\section{Introduction}
Large language models (LLMs) are increasingly used in applications that require long-context processing, including agentic workflows, long-horizon reasoning, and retrieval-augmented generation~\citep{openai2025gpt5, deepmind2025gemini3pro, yang2025qwen3}. As context windows expand, achieving both high throughput and low latency in LLM serving becomes increasingly difficult~\citep{liu2025comprehensive}. This challenge stems from the differing performance characteristics of the two main components of Transformer decoding: the multi-layer perceptron (MLP) and self-attention. The MLP is dominated by model-weight accesses and is typically compute-bound, while self-attention depends on context length and is typically memory-bound~\citep{tang2024quest}.

Batching improves throughput primarily by amortizing model-weight loads across requests in batch, thereby increasing arithmetic intensity in the MLP and shifting its execution from matrix-vector multiplication (GEMV)-like to matrix-matrix multiplication (GEMM)-like computation~\citep{kwon2023efficient, agrawal2024taming, zheng2024sglang}. Attention, however, does not benefit in the same way: each request has an independent Key-Value (KV) cache, so its context cannot be amortized across the batch. For each newly generated token, the model must read the keys and values of the full prior context of that request. Consequently, as sequence length and batch size increase, attention incurs growing memory traffic while remaining low in arithmetic intensity~\citep{agrawal2024vidur}.

To mitigate this memory inefficiency, self-speculative decoding trades computation for lower memory traffic by pairing efficient drafting with exact verification. The model drafts multiple tokens auto-regressively using sparse attention and then verifies them in parallel with a full-attention forward pass. Since the costly access to the full KV cache occurs only once during verification, that cost can be amortized across multiple accepted tokens. As a result, self-speculative decoding is particularly attractive in long-context regimes, where repeated KV-cache reads dominate runtime. Recent systems such as MagicDec~\citep{sadhukhan2025magicdec} and Vegas~\citep{yue2026vegas} demonstrate the effectiveness of this direction.

MagicDec~\citep{sadhukhan2025magicdec} is a batched self-speculative decoding scheme that uses window attention for drafting. Vegas~\citep{yue2026vegas} improves drafting accuracy by selecting high-scoring context tokens based on attention scores from verification tokens. However, Vegas keeps the partial KV context fixed throughout drafting and refreshes it only after verification. As a result, the drafting context becomes increasingly stale as generation progresses, limiting draft length and quality, and ultimately speedup. To address this, we introduce lightweight intra-draft context refresh. By exploiting the similarity of attention patterns across layers for the same token and across consecutive tokens~\citep{zarch-etal-2026-delta, yang2025tidaldecode}, we use a refresh layer that performs full attention during drafting and updates the partial KV for the next token with minimal overhead.

Despite these advances, existing batched speculative decoding systems still operate in a largely synchronized manner: requests in the same batch proceed through drafting and verification together, typically under a shared speculation schedule. While this design simplifies execution, it is poorly matched to real batched workloads. Different requests often have different optimal draft lengths, and even within a single request, the optimal draft length can vary substantially as generation progresses. As a result, synchronized batching forces heterogeneous requests to progress in lockstep, causing some requests to be under-drafted and others to be over-drafted. Thus, the gains of batched speculative decoding are limited not only by suboptimal speculation lengths, but also by synchronization itself.

To address these limitations, we propose ASPIRE, a dynamic, non-synchronized batched speculative decoding framework for LLM inference. ASPIRE is built around a unified mixed batch that allows requests in the same batch to remain in different speculation states within a single model forward pass: some requests draft with partial KV context, while others verify with full KV. This removes the need for synchronized batch-wide draft and verify phases while preserving the efficiency benefits of shared batched execution. On top of this execution model, ASPIRE employs a fine-grained online scheduler that lets each request follow its own speculation trajectory. Using a per-request acceptance-rate estimate and a batch-aware draft/verify cost ratio, the scheduler decides whether each request drafts or verifies in each model batch forward pass. Finally, ASPIRE integrates a lightweight full-attention refresh layer that updates the partial KV context during drafting with minimal overhead, improving draft quality and enabling longer drafting. As a result, ASPIRE improves speculative utilization and boosts throughput for batched long-context LLM inference.

In summary, our contributions are as follows:
\begin{itemize}
    \item We perform an empirical study of draft-length behavior in batched inference and show that optimal draft lengths are highly heterogeneous across requests and highly dynamic over the course of generation within each request.
    \item We propose ASPIRE, a dynamic non-synchronized batched speculative decoding framework that schedules drafting and verification at the granularity of individual requests, allowing requests within the same batch to follow their own speculation trajectories.
    \item We introduce a lightweight context-refresh mechanism during drafting that exploits attention correlation across layers and consecutive tokens to update the partial KV set with minimal overhead.
    \item We conduct extensive experiments across three models and five workloads, showing that ASPIRE substantially improves throughput over AutoRegressive decoding, with up to $4.58\times$ speedup, and consistently achieves the best overall performance against prior speculative decoding baselines.
\end{itemize}
